\section{第六轮：direct epilogue 的低寄存器实验}

\subsection{动机、假设和边界}

第五轮的最小 2SM state2 probe 已经能够 exact 运行，但 NCU 报告
181 registers/thread。该 probe 的 host 端调用固定为
\code{alpha=1,beta=0}，而通用 epilogue 仍然创建了 C 的寄存器 fragment，
执行了一次 C global load 和 AXPBY，再将结果写回 D。因此本轮只检验一个
局部假设：如果把无效的 C/AXPBY 临时量完全移除，寄存器压力是否会下降，
以及下降是否能转化为稳定的 event 加速。

实验不改变 FlashKDA 默认 extension，也不宣称该路径支持任意 alpha/beta。
\code{C1\_TCGEN05\_DIRECT\_EPILOGUE} 只在
\code{C1\_TCGEN05\_STATE2} probe 上启用；通用 AXPBY 分支保留，作为
同一 source、同一 atom、同一输入的对照。该取舍与 CUTLASS Blackwell 文档
列出的 2SM epilogue 调度选项相容，但没有直接把 tutorial epilogue 接入
生产 kernel；参见
\url{https://github.com/NVIDIA/cutlass/blob/main/media/docs/cpp/blackwell_functionality.md}。

\subsection{代码变更}

在 \code{challenge/tcgen05\_sm103\_2sm\_probe.cu} 的 Step 3 中加入
预处理分支：

\begin{lstlisting}
#if defined(C1_TCGEN05_DIRECT_EPILOGUE)
  Tensor tDrAcc = make_tensor<AccType>(shape(tDgD));
  copy(tiled_t2r_copy, tDtAcc, tDrAcc);  // TMEM -> RMEM
  copy(tDrAcc, tDgD);                     // RMEM -> GMEM
#else
  Tensor tDrC = make_fragment_like(tDgC);
  copy(tDgC, tDrC);
  Tensor tDrAcc = make_tensor<AccType>(shape(tDgD));
  copy(tiled_t2r_copy, tDtAcc, tDrAcc);
  axpby(alpha, tDrAcc, beta, tDrC);
  copy(tDrC, tDgD);
#endif
\end{lstlisting}

direct 分支仍使用同一个 TMEM copy atom、同一个输出 partition 和同一个
cluster launch；区别仅在于不生成 \code{tDrC}。这使得寄存器计数差异
可以归因到 epilogue，而不是 tile、TMA 或 launch 配置。

\subsection{B300 构建和运行命令}

源码通过 \code{scp} 同步到 B300 的 challenge 目录，在 Slurm GPU allocation
中以 CUDA 13 的 \code{sm\_103a} 编译。baseline 和 direct binary 使用
同一个 source checkout：

\begin{lstlisting}
cd <B300_REPO>/assignment02/team/c1_flashkda/challenge
nvcc -std=c++17 -O3 -arch=sm_103a --expt-relaxed-constexpr \
  -DC1_TCGEN05_STATE2 -DC1_TCGEN05_QUIET \
  -I../FlashKDA/cutlass/include \
  -I../FlashKDA/cutlass/tools/util/include \
  -I../FlashKDA/cutlass/examples/common -I. \
  -o tcgen05_sm103_2sm_state2_r6_base tcgen05_sm103_2sm_probe_r6.cu

nvcc -std=c++17 -O3 -arch=sm_103a --expt-relaxed-constexpr \
  -DC1_TCGEN05_STATE2 -DC1_TCGEN05_QUIET \
  -DC1_TCGEN05_DIRECT_EPILOGUE \
  -I../FlashKDA/cutlass/include \
  -I../FlashKDA/cutlass/tools/util/include \
  -I../FlashKDA/cutlass/examples/common -I. \
  -o tcgen05_sm103_2sm_state2_r6_direct tcgen05_sm103_2sm_probe_r6.cu
\end{lstlisting}

job 24262 在两个 binary 上分别执行：

\begin{lstlisting}
export C1_TCGEN05_BENCH_ITERS=200 C1_TCGEN05_BENCH_WARMUP=30
./tcgen05_sm103_2sm_state2_r6_base
./tcgen05_sm103_2sm_state2_r6_direct
\end{lstlisting}

两次输出均为 \code{Relative error = 0} 和
\code{Execution is successful}。由于 event 区间只包住已经构造好的
kernel launch，host 端 CuTe layout print 不参与正式时间。

\subsection{结果一：最小 state2 形状}

job 24262 的 \(128\times128\times16\) 对照如下：

\begin{table}[htbp]
\centering
\small
\begin{tabular}{l r r l}
\toprule
variant & event (ms) & registers/thread & correctness \\
\midrule
general AXPBY & 0.0212210 & 181 & exact \\
direct alpha=1,beta=0 & 0.0204546 & 106 & exact \\
\bottomrule
\end{tabular}
\caption{job 24262，200 event iterations、30 次 warm-up。}
\end{table}

direct epilogue 将寄存器数从 181 降到 106，减少
\((181-106)/181=41.4\%\)；CUDA event 从 0.0212210 ms 降到
0.0204546 ms，减少 3.6\%。NCU 单次 replay 的对应时间为 15.584 us
和 14.592 us；该时间受 replay 影响，只用于确认结构变化。

\subsection{结果二：较大 tiled workload}

为了确认收益不是单 tile 噪声，job 24264 使用
\(512\times1024\times64\)（32 个 \(128\times128\times16\) work items）
做一次对照，得到 general 0.0233005 ms、direct 0.0231078 ms；direct
快约 0.8\%。随后 job 24265 在同一 allocation 配置下进行三次成对重复，
每次 300 event iterations、50 次 warm-up：

\begin{table}[htbp]
\centering
\small
\begin{tabular}{c r r r}
\toprule
repeat & general (ms) & direct (ms) & direct speedup \\
\midrule
1 & 0.0174880 & 0.0169668 & 2.98\% \\
2 & 0.0172502 & 0.0169061 & 1.99\% \\
3 & 0.0173580 & 0.0169572 & 2.31\% \\
\bottomrule
\end{tabular}
\caption{job 24265 的三次成对 CUDA-event 重复；两种 variant 的
CPU reference 均 exact。}
\end{table}

重复结果显示 direct 路径有小幅且方向一致的优势，但没有达到本轮预设的
10\% 微内核加速门槛。该门槛用于决定是否值得把路径继续迁移到 K2，
不是作业题目要求的硬性阈值。

\subsection{NCU 资源变化}

job 24263 使用和第五轮相同的 lightweight metric 集合：
\code{launch\_\_block\_\_size}、\code{launch\_\_grid\_\_size}、
\code{launch\_\_registers\_\_per\_\_thread}、
\code{launch\_\_shared\_\_mem\_\_per\_\_block\_\_dynamic}、active warps、
DRAM bytes 和 \code{gpu\_\_time\_\_duration}。两者 launch 均为 block
128、grid (2,1,1)、dynamic shared memory 4,224 B、active warps
约 6.15\%。变化集中在寄存器和 global read：

\begin{table}[htbp]
\centering
\small
\begin{tabular}{l r r r}
\toprule
variant & registers/thread & DRAM read (B) & NCU time (us) \\
\midrule
general AXPBY & 181 & 118,528 & 15.584 \\
direct alpha=1,beta=0 & 106 & 49,664 & 14.592 \\
\bottomrule
\end{tabular}
\caption{job 24263 的 NCU CSV；DRAM write 计数在该轻量 replay 中为零，
不能解释为没有 global store。}
\end{table}

direct 路径减少了 C tile 的读取，因此 DRAM read 也明显下降；但对更大
问题，固定的 TMEM allocation、2SM barrier、TMA multicast 和 output store
仍占据主要时间。Blackwell tuning guide 将寄存器文件、shared memory 和
cluster occupancy 都列为需要同时检查的资源，参见
\url{https://docs.nvidia.com/cuda/blackwell-tuning-guide/}。

\begin{record}{R6 结论与接入决策}
本轮实现了一个可回滚的 opt-in direct epilogue，并在 B300 上证明：
\begin{enumerate}
\item alpha=1,beta=0 state2 probe 保持 exact；
\item registers/thread 从 181 降到 106，证明通用 AXPBY 临时量确实是
      当前微内核寄存器压力的重要来源；
\item 最小 tile event 快 3.6\%，较大 workload 的三次重复快
      1.99--2.98\%，但没有达到 10\% 的迁移门槛；
\item 该优化不能直接用于完整 K2，因为 K2 的 state update、alpha/beta
      语义和输出布局尚未与这个 GEMM probe 对齐。
\end{enumerate}

因此第六轮不把 direct 分支接入 \FlashKDA 默认路径，也不把它宣称为端到端
baseline 超越。下一轮若继续，优先顺序应为：先构造保留 K2 state 语义的
state-only 写回 probe；只有在该 probe 同时保持 exact 且出现稳定的双位数
收益时，才值得进行一次窄范围 K2 集成。否则应保留 baseline，并将本轮的
41.4\% register reduction 作为 SM100 v2 的资源设计证据。
\end{record}
